\section{The Coming}

\subsection{The Ox Knows His Owner}\label{sec:ox}

On the first night of Advent the Church opens the book of Isaiah. At Matins of
the First Sunday the Breviary begins the prophet from his first verse. Isaiah
sees his vision, calls on heaven and earth to listen, and then says this:

\begin{quote}
The ox knoweth his owner, and the ass his master's crib: but Israel hath not
known me, and my people hath not understood.\par
{\smallskip\noindent\small ---Isaiah 1:3\par}
\end{quote}

The crib is a manger. Where the Douay has ``his master's crib'' the Latin of the
lesson reads \latin{praesepe domini sui}, and \latin{praesepe} is the word of the
Latin Gospel when St Luke tells how Mary laid her Son \latin{in praesepio}, in a
manger. The season's first lesson speaks of a beast standing at the trough
where it is fed, and of a people who do not know the hand that feeds them.

Four weeks later the Church comes back to the same trough. At Matins of
Christmas, when the first lesson from St Leo has been read, she sings the
responsory that follows it: \latin{O magnum mysterium, et admirabile
sacramentum, ut animalia viderent Dominum natum, jacentem in praesepio}. It
calls the birth a great mystery and a wonderful sacrament, and it sings of the
manger: ``The beasts saw the new-born Lord lying in a manger.'' The responsory
does not say which beasts.

St Gregory Nazianzen does. Preaching on the day of the Nativity, he
takes his hearers through the whole mystery, the census, the birth, the little
town and the crib, and then hands Isaiah's verse to them as a command:

\begin{quote}
Revere the enrolment on account of which thou wast written in heaven, and adore
the Birth by which thou wast loosed from the chains of thy birth, and honour
little Bethlehem, which hath led thee back to Paradise; and worship the manger
through which thou, being without sense, wast fed by the Word. Know as Isaiah
bids thee, thine Owner, like the ox, and like the ass thy Master's crib; if thou
be one of those who are pure and lawful food, and who chew the cud of the word
and are fit for sacrifice.\par
{\smallskip\noindent\small ---St Gregory Nazianzen, Oration 38, on the Birthday of Christ\par}
\end{quote}

St Ambrose, writing on St Luke's account of the birth, hears the same animal.
He writes, \latin{infantis audis vagitus, non audis bovis dominum agnoscentis
mugitus}: you hear the cry of the infant, you do not hear the lowing of the ox
that knows its Lord. And he goes on at once to Isaiah, \latin{agnovit enim bos
possessorem suum et asinus praesaepe domini sui}: for the ox knew his owner,
and the ass his master's crib.

On the first night of Advent the Church reads this against herself. The ox
does not reason about its owner; it knows him, and it knows where it is fed.
Israel, made for God, did not know him. Gregory turns the reproach into an
invitation. We were the creatures ``without sense'', and at the manger we were
``fed by the Word''. The trough at Bethlehem is the place where the Owner whom
we did not know makes himself known, and where those he made are fed.

Advent, then, is four weeks of learning to know the One who comes. The Church
does not begin with the crib. She begins on the First Sunday with the Gospel of
the end of the world, and she ends on the Vigil, at the door of Bethlehem.
Between the two lie four Sundays; the Ember Wednesday, Friday and Saturday of
the third week, days of fasting that keep the old December fast of the Roman
Church; and the last days from 17 December, when Vespers is crowned each
evening with one of the great O antiphons.

The book of Isaiah runs beside the whole season. Day by day the Office reads
it from its first chapter to its last, and the Church comes to the manger with
the prophet's words in her mouth.

One question comes first. The Church says she is waiting for the Lord's
coming, and she has already received him. Whose coming does she keep, and how
many comings are there?
